\section{Method}
\label{sec:method}

\paragraph{Localization task and evaluation denominators}
EAPG takes one evaluated C++17 program and a size-parameterized workload, and ranks canonical source locations in that program. A \emph{work symptom} is a node whose measured operation-weighted growth satisfies an alert rule; the \emph{reference} marks where a single-site edit alters workload growth, such as resetting a monotonic scan cursor inside an outer loop. Fast or restored versions and edit references stay outside prediction. In the archived pilot, a mechanism-informed AST diff designates one reference node per pair. In the prospective external design, \texttt{score\_external.py} joins predictions against sealed ground truth over drawn pairs using three nested positive denominators: $D_{\mathrm{pos}}$ (all resolved confirmed positives, counting abstentions or normalization failures as misses), $D_{\mathrm{elig}} \subseteq D_{\mathrm{pos}}$ (canonical node count $|V|\le 26$ and serialized prompt length $+1 \le 512$ tokens), and $D_{\mathrm{rank}} \subseteq D_{\mathrm{elig}}$ (slices with $|C|\ge 2$), alongside false alerts on verified negatives and conservative lower bounds counting unresolved positives as misses. The prospective primary estimand is the equal-family mean of paired Top-1 differences between scoped and conventional slices under identical work symptoms and ranking policies.

\paragraph{Canonical representation and AST def/use approximation}
Pass~0 strips literal \texttt{\#include} directives, expands remaining macros via \texttt{g++ -E -P}, and restores the directives. A tree-sitter-cpp pass inserts braces around supported single-statement loop and branch bodies in reverse byte order and validates syntax with \texttt{g++ -std=c++17 -fsyntax-only}. A pre-order walk assigns canonical node IDs (\texttt{N1}, \texttt{N2}, \dots) to function entries, loops, branches, calls, and state updates, partitioned by an uninstrumented \texttt{MAIN\_SCOPE} boundary. Each macro-node records its parent scope, byte span, and approximate variable sets $\mathrm{defs}(u)$ and $\mathrm{uses}(u)$. Explicitly, for any \texttt{for\_statement} with an update expression (e.g., \texttt{b += jump}), the extractor places all identifiers parsed inside the update expression into both $\mathrm{defs}(u)$ and $\mathrm{uses}(u)$, over-approximating read-only step variables as loop definitions. The prospective candidate (\texttt{2c95df0e}) refines receiver type bindings to respect lexical scopes and limits inferred caller edges to direct calls targeting visible free functions.

\paragraph{Work proxies, tick ceiling, and scaling fits}
Each instrumented event increments hit count $h_v(N)$ by 1 and work count $w_v(N)$ by a static weight expression. Ticks fire on entry to function/loop bodies, before branches and calls, and after state updates; the prospective revision also places ticks inside supported call expressions in conditions and assignments so short-circuited calls contribute no tick. Unit weight defaults to $1$, whereas qualified \texttt{std::find}/\texttt{std::count} calls (plus \texttt{std::reverse} in the prospective revision) and recognized container member calls receive linear size $s$ or tree $\lfloor\log_2(s+2)\rfloor$ proxies, with hash-container lookups weighted by $1$. Crucially, the archived pilot (\texttt{pipeline.py}) defaulted unknown container receivers to linear size, whereas the prospective revision defaults unknown receivers to $1$. Binaries compile with GCC \texttt{-O2} under a dynamic ceiling
\begin{equation}
B(N) = \max\!\left(2\times 10^6,\; 5N^2\right),
\end{equation}
using intermediate ladder probes below the first noncompletion. Hitting $B(N)$ emits a \texttt{CAP\_BREACHED} JSON record via \texttt{std::exit(0)}, so process exit code~$0$ does not imply clean completion; only uncapped completions enter fits. For node $v$ and channel $c \in \{\mathrm{hits},\mathrm{work}\}$, OLS estimates power-law ($\log\max(1,Y_v^c(N))$ vs.\ $\log N$, slope $\alpha_v^c$) and exponential ($\log_2\max(1,Y_v^c(N))$ vs.\ $N$, slope $\beta_v^c \ge 0.02$ over $\ge 3$ points with strictly higher $R^2 \ge 0.85$) growth across uncapped sizes. Prospective comparison requires four clean sizes spanning at least eightfold, an active series, $R^2 \ge 0.85$ for the selected family, independent operation counts, output checks, and exact restoration. Separately from counted operations, uninstrumented timing confirmation requires a log-log slope of median slow-to-fast wall time $\ge 0.50$ on every input seed and a largest-size ratio $\ge 1.50$.

\paragraph{Symptom selection and scoped backward slice}
A node $s$ qualifies as a work symptom under three triggers: (i) an exponential onset relative to parent scope, (ii) a hidden-work gap $\alpha_s^{\mathrm{work}} - \alpha_s^{\mathrm{hits}} \ge 0.70$ (requiring no floor on $\alpha_s^{\mathrm{work}}$), or (iii) a scope-relative work jump $\Delta \alpha_s = \alpha_s^{\mathrm{work}} - \alpha_{\mathrm{parent}(s)}^{\mathrm{work}} \ge 0.70$. Whereas the archived selector used traversal order and fell back conservatively, the repaired prospective rule admits top-level loops under \texttt{MAIN\_SCOPE} (where $\Delta\alpha_s = \alpha_s^{\mathrm{work}}$), enforces $\alpha_s^{\mathrm{work}} \ge 1.50$ on scope-relative jumps (not hidden-work calls), and orders exponential fits before larger polynomial exponents. Given symptom $s$ with guard variables $V_{\mathrm{guard}}(s) = \mathrm{uses}(s)$ (augmented by parameters and caller uses for functions), the pruned scoped slice is $C_{\mathrm{pruned}}(s) = \{s\} \cup C_{\mathrm{inter}}(s) \cup C_{\mathrm{intra}}(s)$, selecting only:
\begin{align}
C_{\mathrm{inter}}(s) &= \bigl\{u : \mathrm{par}(u)=\mathrm{par}(s),\; \mathrm{defs}(u)\cap V_{\mathrm{guard}} \neq \emptyset, \nonumber\\
&\qquad (\mathrm{exp}(s) \lor \alpha_u^{\mathrm{work}} \le \alpha_s^{\mathrm{work}} - 0.50)\bigr\},\\
C_{\mathrm{intra}}(s) &= \bigl\{u : \mathrm{par}(u)=s,\; \mathrm{defs}(u)\cap V_{\mathrm{guard}} \neq \emptyset, \nonumber\\
&\qquad (\mathrm{reset}(u) \lor \mathrm{nested}(u)) \land \neg\mathrm{step}(u)\bigr\}.
\end{align}
Because $C_{\mathrm{inter}}$ inspects only $\mathrm{parent}(u) = \mathrm{parent}(s)$ (bypassing the $-0.50$ exponent guard when $s$ is exponential), an edit two scopes above $s$ or connected solely through intermediate copies lies outside $C_{\mathrm{pruned}}(s)$, while unkilled low-growth writes to $V_{\mathrm{guard}}$ can enter it. Candidate space reduction is $R_{\mathrm{search}} = 1 - |C_{\mathrm{pruned}}(s)|/|V|$.

\paragraph{Comparators, six-representation matrix, and contract}
The conventional comparator constructs reaching definitions and control dependences on a structured CFG for single-\texttt{main()} integral scalar C++17 programs~\cite{weiser1981slicing}, projecting backward closure from $s$ onto canonical nodes. Within fixed candidate slices, we compare four non-neural policies: uniform random guessing ($\mathbf{E}[\mathrm{Top\text{-}1}]=1/|C|$), raw AST traversal order, the prospective default symbolic ranker (\nolinkurl{use_loop_def_use=False}), and the pilot-refined symbolic rule (\nolinkurl{use_loop_def_use=True}) that filters loop-bound aliases using approximated $\mathrm{defs}$ sets. Crucially, \nolinkurl{rank_candidates_by_role_priority} unconditionally hoists all non-symptom \texttt{STATE\_UPD} nodes surviving inside $C \setminus \{s\}$ (intra-resets before inter-resets) ahead of $s$, unless an inner flag-without-break pattern preserves $s$ first; hence any CFG slice carrying initializers or loop accumulators into \texttt{STATE\_UPD} ranks them before $s$. To isolate which signals drive neural localization, we evaluate 4-bit NF4 \texttt{Qwen2.5-Coder-1.5B}~\cite{hui2024qwen,dettmers2023qlora} (base and QLoRA-adapted) across six single-token-classification prompt serializations capped at $\le 512$ tokens: (1) \nolinkurl{raw_code} (26-char node snippets), (2) \nolinkurl{hotspot_raw} (code snippets plus hit exponents $\alpha_v^{\mathrm{hits}}$), (3) \nolinkurl{static_eapg} (roles, identifier $\mathrm{defs}/\mathrm{uses}$, code snippets), (4) \nolinkurl{abstract_hits_only} (roles, D/U counts, $\alpha_v^{\mathrm{hits}}$, binary \texttt{SLICE} tags, omitting $\alpha_v^{\mathrm{work}}$), (5) \nolinkurl{abstract_graph} (binary-tagged roles, D/U counts, both $\alpha_v^{\mathrm{hits}}/\alpha_v^{\mathrm{work}}$, \texttt{SLICE} tags), and (6) \nolinkurl{full_eapg} (\nolinkurl{abstract_graph} with D/U counts replaced by truncated lists of up to two def/use identifiers plus 26-char code snippets). One forward pass yields logits over all canonical nodes \texttt{A..Z}, permitting paired \emph{Unconstrained} ($V$) and \emph{Slice-Constrained} ($C_{\mathrm{pruned}}$) decoding. For prospective deployment, \nolinkurl{predict_external_contract.py} enforces identical tokenization and single-token checks (\nolinkurl{prediction_contract.py}) bound by adapter/protocol SHA-256 digests and emits explicit abstention codes (\nolinkurl{NO_SYMPTOM}, \nolinkurl{EMPTY_SLICE}, \nolinkurl{NODE_LIMIT}, \nolinkurl{CONTEXT_LIMIT}, \nolinkurl{UNSUPPORTED_NODE_TOKEN}, \nolinkurl{NORMALIZATION_FAILURE}) without falling back to $V$.
